\chapter{Οδηγός Εγκατάστασης και Χρήσης}
\label{app:usage}

\section{Απαιτήσεις Συστήματος}

Για την εκτέλεση του project απαιτείται:
\begin{itemize}
    \item Python 3.10 ή νεότερο
    \item \texttt{pip} για εγκατάσταση εξαρτήσεων
    \item βασικές βιβλιοθήκες επιστημονικού Python stack:
    \texttt{numpy}, \texttt{scipy}, \texttt{pandas}, \texttt{matplotlib}
    \item προαιρετικά: \texttt{streamlit} για το Python UI
    \item προαιρετικά: \texttt{node/npm} για τη Next.js εφαρμογή
\end{itemize}

Η τρέχουσα έκδοση του αποθετηρίου έχει επαληθευθεί με \texttt{Python 3.12.2}. Για πλήρη έλεγχο του περιβάλλοντος διατίθεται το script \filepath{verify_setup.py}.

\section{Εγκατάσταση}

\subsection{Βασική Εγκατάσταση}

\begin{lstlisting}
python3 -m venv venv
source venv/bin/activate
pip install -r requirements.txt
\end{lstlisting}

\subsection{Έλεγχος Περιβάλλοντος}

\begin{lstlisting}
python verify_setup.py
\end{lstlisting}

\subsection{Εκτέλεση Test Suite}

\begin{lstlisting}
python -m pytest tests -q
\end{lstlisting}

\section{Βασική Χρήση}

Η απλούστερη χρήση του project είναι μέσω των demo scripts. Τα scripts αυτά καλύπτουν δημιουργία κύβου, scrambling, επίλυση και σύγκριση αλγορίθμων.

\subsection{Εκτέλεση Demos}
\begin{lstlisting}
python demos/basic_usage.py
python demos/kociemba_demo.py
python demos/thistlethwaite_demo.py
python demos/a_star_comparison_demo.py
python demos/distance_estimator_demo.py
\end{lstlisting}

\subsection{Benchmarks και Συγκρίσεις}

\begin{lstlisting}
python demos/phase9/benchmark_demo.py --n-scrambles 20
python demos/phase9/algorithm_comparison_cli.py --depth 10 --seed 42
\end{lstlisting}

\section{Διαδραστικές Εφαρμογές}

\subsection{Streamlit}

\begin{lstlisting}
streamlit run ui/app.py
\end{lstlisting}

\subsection{Next.js Web Application}

\begin{lstlisting}
cd webapp
npm ci
npm run build
npm run dev
\end{lstlisting}

Η εφαρμογή Next.js λειτουργεί ως presentation/demo frontend και όχι ως authoritative live διεπαφή εκτέλεσης των Python solvers. Παρέχει σελίδες για:
\begin{itemize}
    \item αρχική παρουσίαση του project,
    \item demo επίλυσης μεμονωμένου scramble,
    \item demo σύγκρισης αλγορίθμων,
    \item και εκπαιδευτικό περιεχόμενο.
\end{itemize}
Η ζωντανή εκτέλεση των αλγορίθμων γίνεται από το \filepath{ui/} και από τα scripts benchmark/validation του repository.

\section{Αναπαραγωγή Αποτελεσμάτων Διπλωματικής}

Για την ακριβή αναπαραγωγή των αποτελεσμάτων του Κεφαλαίου~\ref{ch:evaluation}
χρησιμοποιείται ως επίσημη διαδρομή το script
\filepath{scripts/benchmarks/regenerate_thesis_benchmarks.py}, το οποίο
επαναεκτελεί τα benchmarks πάνω στο αποθηκευμένο σύνολο scrambles και
αναδημιουργεί τα αρχεία
\filepath{results/benchmarks/thesis/thesis_bench_d*.json} και
\filepath{results/benchmarks/thesis/thesis_results_combined.json}. Προαιρετικά
μπορούν να εκτελεστούν και τα scripts ανάλυσης, τα οποία διαβάζουν τα παραπάνω
αρχεία από τον ίδιο κατάλογο:

\begin{lstlisting}
python scripts/benchmarks/regenerate_thesis_benchmarks.py
python scripts/benchmarks/analyze_thesis_data.py
python scripts/benchmarks/generate_latex_tables.py
\end{lstlisting}

Τα scripts \filepath{generate_thesis_data.py} και
\filepath{generate_complete_thesis_data.py} διατηρούνται μόνο για ιστορικούς ή
διερευνητικούς λόγους και δεν αποτελούν τη διαδρομή αναπαραγωγής της
διπλωματικής.

Για τη διαχείριση της συγγραφής και των chapter packets προστέθηκε επίσης workflow script:

\begin{lstlisting}
python scripts/thesis_workflow.py status
python scripts/thesis_workflow.py packets --remaining
python scripts/thesis_workflow.py validate
\end{lstlisting}

\section{Σημειώσεις}

\begin{itemize}
    \item Η πρώτη εκτέλεση ορισμένων solvers μπορεί να είναι αισθητά πιο αργή λόγω δημιουργίας ή φόρτωσης πινάκων.
    \item Τα benchmark scripts αποθηκεύουν ρητά τα scrambles και τα αποτελέσματα σε JSON μορφή.
    \item Η πλήρης παραγωγή PDF της διπλωματικής απαιτεί TeX toolchain συμβατό με το \filepath{thesis/main.tex}. Στο αποθετήριο χρησιμοποιείται \texttt{tectonic}.
\end{itemize}
